% language=us runpath=texruns:manuals/luametafun

\environment luametafun-style

\startcomponent luametafun-voronoi

\startchapter[title={Voronoi and Delaunay}]

This is a world of its own and we can't cover everything into this chapter so
examples have to be the introduction. It might help to search the Internet first
in order to understand that these graphics and related methods are widely used.
It might also become clear that we have a computational heavy mechanism here,
although for average use performance should be okay.

Before we continue it is important to stress that we don't actually extend
\METAPOST\ with such features. It is way more natural to do these things in \LUA,
maybe with some help from small \CCODE\ we provide, like efficiently handling
coordinates, so that we keep the \METAPOST\ code base clean and focused. As with
other extensions, like 3D surfaces, we try to provide a simple and consistent
key/value interface to the \LUA\ end. Also, if you really want more fancy things,
you should invest time in the often very rich graphical user interfaced
applications. There is no way we can and will compete with that, if only because
there is no place for evolving libraries in what is basically a core typesetting
application. Also, we want user interfaces to be consistent, familiar and
somewhat predictable in \CONTEXT.

\startbuffer
\startMPcode
    path p ; p := lmt_voronoi [
        solution = "voronoi",
        method   = "path",
        name     = "demo",
        width    = 200,
        height   = 100,
        points   = 250,
        seed     = 0.123,
    ] xsized TextWidth ;
    for i = 1 upto segments p :
        fill (subpath (segment i of p) of p) && cycle
            withcolor (.4 randomized .2,.4 randomized .2,.4 randomized .2)
        ;
    endfor ;
\stopMPcode
\stopbuffer

\startplacefigure[reference=fig:voronoi:1,title=Getting the result as path]
    \getbuffer
\stopplacefigure

% \stoptext

This graphic is generated with the following code. You can provide a dataset as we will
see later, but in this case we generate a random set of points.

\typebuffer[option=TEX]

\in {Figure} [fig:voronoi:2] show a similar graphic but here save the result and
loop over the cells but we actually have better ways to do that using helpers.
Here we show \emphasize {cells} but there are also \emphasize {triangles},
\emphasize {points} and \emphasize {edges}.

\startbuffer
\startMPcode
    lmt_voronoi [
        solution = "voronoi",
        method   = "store",
        name     = "demo",
        width    = 200,
        height   = 100,
        points   = 250,
        seed     = 0.456,
    ] ;

    draw image (
        for i = 1 upto stpcount "demo" "cells" :
            fill (stpsegment "demo" "cells" i) && cycle
                withcolor (.4 randomized .2)
            ;
        endfor ;
        draw (stppath "demo" "cells") && cycle
            withpen pencircle scaled .5mm
            withcolor "darkgreen" ;
        for i = 1 upto stpcount "demo" "cells" :
            draw textext ("\infofont " & decimal i) scaled .25
                shifted centerofmass (stpsegment "demo" "cells" i)
                withcolor "white"
            ;
        endfor ;
    ) xsized TextWidth ;
\stopMPcode
\stopbuffer

\startplacefigure[reference=fig:voronoi:2,title=Rendering the stored result explicitly.]
    \getbuffer
\stopplacefigure

In the next example we use a predefined set of points. Watch how we use the vector
library! The generators use these efficient point lists.

\startbuffer
\startluacode
 -- local random = math.random
    local random = utilities.randomizer.get

    local n    = 40
    local nn   = n * n // 2
    local set  = vector.points.new(nn,1,2)
    local tree = table.setmetatableindex("table")

    set(1,1) tree[1][1] = true
    set(1,n) tree[1][n] = true
    set(n,n) tree[n][n] = true
    set(n,1) tree[n][1] = true

    for i=1,nn do
        while true do
            local x = random(1,n)
            local y = random(1,n)
            if not tree[x][y] then
                set(x,y)
                tree[x][y] = true
                break
            end
        end
    end

    metapost.registerpointset {
        name   = "myset-3",
        points = set
    }
\stopluacode
\stopbuffer

\typebuffer[option=TEX]

These points are used in the rendering of \in {figure} [fig:voronoi:3] where we
show the four possible results. The dimensions are important because we need an
outer edge to be able to \quote {clip} or \quote {bound} the result. Internally
a temporary huge triangle is used to help with this (commonly used trickery).

\startbuffer
\startMPcode
    draw image (
        lmt_voronoi [
            solution = "voronoi",
            method   = "store",
            name     = "context",
            pointset = "myset-3",
            width    = 41,
            height   = 41,
        ] ;

        draw (stppath "context" "cells")
            withpen pencircle scaled .25
            withcolor "darkgreen"
        ;
        draw (stppath "context" "edges")
            xshifted 50
            withpen pencircle scaled .25
            withcolor "darkblue"
        ;
        draw (stppath "context" "triangles")
            xshifted 100
            withpen pencircle scaled .25
            withcolor "darkred"
        ;
        drawdot (stppath "context" "points")
          % xshifted 150
            withpen pencircle scaled 1
            withcolor "darkyellow"
        ;
    ) xsized TextWidth ;
\stopMPcode
\stopbuffer

\typebuffer[option=TEX]

\startplacefigure[reference=fig:voronoi:3,title=Four renderings using given points.]
    \getbuffer
\stopplacefigure

The Voronoi renderer first does a Delaunay that collects te edges, neighbors and
triangles, followed by cell construction. In \in {figure} [fig:voronoi:4] we see
what we get when we ask for various paths.

\startbuffer
\startMPcode
    draw image (
        lmt_voronoi [
            solution = "delaunay",
            method   = "store",
            name     = "context",
            pointset = "myset-3",
            width    = 41,
            height   = 41,
        ] ;

        draw (stppath "context" "cells")
            withpen pencircle scaled .25
            withcolor "darkgreen"
        ;
        draw (stppath "context" "edges")
            xshifted 50
            withpen pencircle scaled .25
            withcolor "darkblue"
        ;
        draw (stppath "context" "triangles")
            xshifted 100
            withpen pencircle scaled .25
            withcolor "darkred"
        ;
        drawdot (stppath "context" "points")
          % xshifted 150
            withpen pencircle scaled 1
            withcolor "darkyellow"
        ;
    ) xsized TextWidth ;
\stopMPcode
\stopbuffer

\typebuffer[option=TEX]

\startplacefigure[reference=fig:voronoi:4,title=Delaunay is just Voronoi without cells.]
    \getbuffer
\stopplacefigure

This whole sub-project started actually by the wish to have stippled text and
what better to use that the word \CONTEXT. The text is typeset with:

\typefile[option=TEX]{luametafun-voronoi-001.tex}

Because here we want to show a few variants we define a graphic. We can pass a
bytemap number of a filename. Three variants are shown in \in {figure}
[fig:voronoi:5].

\startbuffer
\startuseMPgraphic{context}{iterations}
  % triplet bm ;
  % bm := loadbytemapfromfile(5, "luametafun-voronoi-001.png") ;

    drawdot lmt_voronoi [
        solution     = "stipple",
        filename     = "luametafun-voronoi-001.png",
        bytemap      = 5,
        name         = "context",
        points       = "auto",
        densityscale = .075,
        iterations   = mpvarn "iterations",
        seed         = 0.5,
        gamma        = 1.15,
        contrast     = 1.08,
        step         = 1,
    ]
        xsized TextWidth
        withpen pencircle scaled 1
        withcolor "white" ;
    addbackground
        withcolor "darkblue" ;
\stopuseMPgraphic
\stopbuffer

\getbuffer \typebuffer[option=TEX]

The parameters \type {gamma}, \type {contrast}, \type {invert}, \type {threshold}, \type
{minimumdensity}, \type {maximumdensity} and \type {densityscale} determine how
the gray scale or luminance values taken from the \RGB\ components are converted.

A \type {value} at a specific point is transformed as follows:

\starttyping[option=LUA]
if invert then
    value = 1 - value
end
if contrast then
    value = (value - 0.5) * contrast + 0.5
    if value < 0 then
        value = 0
    elseif value > 1 then
        value = 1
    end
end
if threshold then
    if value <= threshold then
        value = 0
    elseif threshold < 1 then
        value = (value - threshold) / (1 - threshold)
    end
end
if gamma then
    value = value ^ gamma
end
value = minimumdensity + (maximumdensity - minimumdensity) * value
\stoptyping

just that you know it. At some point we might come up with recommendations.

\startplacefigure[reference=fig:voronoi:5,title=Stippling quality is determined by iterations.]
    \startcombination[nx=1,ny=3]
        {\useMPgraphic{context}{iterations=0}}   {0 iterations}
        {\useMPgraphic{context}{iterations=4}}   {4 iterations}
        {\useMPgraphic{context}{iterations=16}} {16 iterations}
    \stopcombination
\stopplacefigure

In this example you should notice that we use \type {drawdot} on the returned
path. This is more efficient than looping over the points in the path and drawing
each point separately. In the next example we use of course a regular \type
{draw}, resulting in \in {figure} [fig:voronoi:6]. When generating this chapter,
rendering the image takes longer than generating these paths.

\startbuffer
\startuseMPgraphic{context}{detail,color}
    draw lmt_voronoi [
        solution     = "stipple",
        filename     = "luametafun-voronoi-001.png",
        bytemap      = 5,
        name         = "context",
        points       = "auto",
        densityscale = .075,
        iterations   = 16,
        seed         = 0.5,
        gamma        = 1.15,
        contrast     = 1.08,
        step         = 1,
        detail       = mpvars "detail"
    ]
        xsized TextWidth
        withpen pencircle scaled .5
        withcolor mpvars "color" ;        ;
\stopuseMPgraphic
\stopbuffer

\getbuffer \typebuffer[option=TEX]

\startplacefigure[reference=fig:voronoi:6,title=You can also get the triangles and cells.]
    \startcombination[nx=1,ny=2]
        {\useMPgraphic{context}{detail=triangles,color=darkgreen}} {triangles}
        {\useMPgraphic{context}{detail=cells,color=darkblue}}      {cells}
    \stopcombination
\stopplacefigure

Next we show some randomizing. Let's pretend that we immediately see which one is
a poisson distribution; see \in {figure} [fig:voronoi:7]. The code also demonstrates
how we can use loops to get six figures from one specification.

\startbuffer[a]
\startMPcode
draw image (
    numeric dy ; dy := 0;
    for how = "poisson", "uniform" :
        numeric dx ; dx := 0;
        for what = "points", "cells", "triangles" :
            if what == "points" : drawdot else : draw fi
            lmt_voronoi [
                solution    = "cvd",
                name        = "context",
                initializer = how,
                points      = 200,
                iterations  = 20,
                seed        = .5,
                detail      = what,
            ]
                shifted (dx,-dy)
                withcolor if how == "poisson" : "darkred" else : "darkblue" fi
                withpen pencircle scaled 2
            ;
            dx := dx + 110 ;
        endfor ;
        dy := dy + 110 ;
    endfor ;
) xsized TextWidth ;
\stopMPcode
\stopbuffer

\startbuffer[b]
\startMPcode
draw image (
    for how = "poisson", "uniform" :
        for what = "points", "cells", "triangles" :
            if what == "points" : drawdot else : draw fi
            lmt_voronoi [
                solution    = "cvd",
                name        = "context",
                initializer = how,
                points      = 200,
                iterations  = 20,
                seed        = .5,
                detail      = what,
            ]
                shifted (
                    110 * (      iteratorindex    - 1),
                    110 * ((deltaiteratorindex 1) - 1)
                )
                withcolor if odd(deltaiteratorindex 1) : "darkred" else : "darkblue" fi
                withpen pencircle scaled 2
            ;
        endfor ;
    endfor ;
) xsized TextWidth ;
\stopMPcode
\stopbuffer

\typebuffer[a][option=TEX]

\startplacefigure[reference=fig:voronoi:7,title=Poisson versus uniform randomization.]
    \getbuffer[b]
\stopplacefigure

Let's use this opportunity to educate you on loops. Iterating over a list of
expressions (here strings) is a natural \METAPOST\ features but we added two
accessors to the current position in these lists. That way we can avoid the
temporary \type {dx} and \type {dy} variables and toggle a color on an index
being odd.

\typebuffer[b][option=TEX]

So how do the methods differ? The \type {delaunay} is the simplest, it generated
edges and triangles. With \type {voronoi} we also get (often non-triangular)
cells. The \type {weighted} and \type {unweighted} distribute points so they do
even more. In this stage of development we permit ourselves to be a bit vague
and stick to examples. Feel free to play around and feedback.

\stopchapter

\stopcomponent

% timestamp : Snarky Puppy + Metropole, North Sea Jazz 2026 (NTR: TV) https://youtu.be/kgYfBQs5_tI
